\begin{abstract}
A few measured frequency-response functions determine a second-order model
of a structure in modal coordinates: its poles, its rank-one residues and
the rows of the modal matrix at the measured ports. For this fitted model
the modal displacement at every frequency follows from the measured forces
alone, however few the sensors. The measurements leave open only how the
modes are distributed over the unmeasured degrees of freedom. A nominal
finite-element model closes this gap. The completed modal matrix keeps the
measured ports exactly and keeps all modal coordinates on the mesh. Among
all such matrices it minimizes the residual of the nominal equation of
motion along the response of the fitted model, with its identified damping,
to the measured forces. The nominal matrices themselves are never fitted.
The completed model is modal on the whole mesh, and its receptance answers
queries: complex frequency responses between known load and response
locations, or forces at known positions from measured accelerations. A
three-mass chain makes every step explicit. A laboratory shell shows exact
preservation of the measured transfers, the limits of spatial prediction
under an imperfect prior, and force estimation with one and two active
shakers.
\end{abstract}
